%! Characteristics of Linear Functions — Unit 1, Lesson 1.0 — Guided Notes
% Gradual release (user direction, 2026-08-31): the notes are the "I do / we do" block.
% 15 minutes. Vocabulary is built here, BEFORE the group activity uses it.
\documentclass[10pt]{article}
\usepackage{algebra2-article}
\usepackage{algebra2-boxes}
\newcommand{\TallMath}[1]{$\displaystyle #1\rule[-1.4em]{0pt}{3.2em}$}
% Coordinate grid: {xmin}{xmax}{ymin}{ymax}
\newcommand{\lgrid}[4]{%
  \draw[step=1, linegray!40, very thin] (#1,#3) grid (#2,#4);
  \draw[->, semithick, charcoal] (#1,0)--(#2,0) node[below right, font=\scriptsize]{$x$};
  \draw[->, semithick, charcoal] (0,#3)--(0,#4) node[left, font=\scriptsize]{$y$};}

\begin{document}

\pageheader{Unit 1, Lesson 1.0}{Guided Notes}

\vspace{0.10in}

\begin{objectivebox}
\textbf{By the end of this lesson, I will be able to\dots}
\begin{itemize}\setlength{\itemsep}{2pt}
  \item recognise a \emph{linear} function by its \emph{constant rate of change}, and find that rate
        (the \emph{slope}) from a graph, a table, or two points;
  \item find the \emph{$y$-intercept} and the \emph{zero} ($x$-intercept) from a graph, a table, or a rule;
  \item say where a line is \emph{increasing}, \emph{decreasing}, or \emph{constant}, and where it is
        \emph{positive} or \emph{negative};
  \item state the \emph{domain} and \emph{range} of a line --- as pure math, and as a story allows.
\end{itemize}
\end{objectivebox}

\vspace{0.05in}

\begin{vocabbox}
\small
Fill in each term as we build it together.
\par\vspace{2pt}   % \par is required: \termblanklong's \noindent is a no-op mid-paragraph
\termblanklong{Linear function}
\termblanklong{Slope (rate of change)}
\termblanklong{$y$-intercept}
\termblanklong{$x$-intercept (zero)}
\termblanklong{Increasing / decreasing / constant}
\termblanklong{Positive / negative interval}
\end{vocabbox}

\begin{hookbox}
At 6:00 a.m.\ Saturday it was $-6\,^\circ$C in Blacksburg, and the temperature rose a steady
$2\,^\circ$C every hour. Let $h$ be the hours after 6:00 a.m.
\begin{itemize}\setlength{\itemsep}{1pt}
  \item $T(0)=\blank{1.0cm}$ \quad $T(1)=\blank{1.0cm}$ \quad $T(2)=\blank{1.0cm}$ \quad $T(4)=\blank{1.0cm}$
  \item The temperature reaches $0\,^\circ$ at $h=\blank{1.0cm}$ --- and \emph{that} is the hour the
        roads stop being icy.
\end{itemize}
Two numbers tell the whole story: \emph{where it starts} and \emph{how fast it changes}. Today we
name every feature of a line, and read each one three ways --- off a graph, a table, and a rule.
\end{hookbox}

\boxguard[24]
\begin{notesbox}{1. Linear Means the Rate of Change Is Constant}
Fill in the Saturday table, then the row of changes underneath it.
\begin{center}
\renewcommand{\arraystretch}{1.3}
\begin{tabular}{|c|c|c|c|c|c|}
\hline
$h$ & $0$ & $1$ & $2$ & $3$ & $4$ \\ \hline
$T$ & $-6$ & $-4$ & \blank{0.9cm} & \blank{0.9cm} & \blank{0.9cm} \\ \hline
change in $T$ & --- & $+2$ & \blank{0.9cm} & \blank{0.9cm} & \blank{0.9cm} \\ \hline
\end{tabular}
\end{center}
The change is the \blank{2.4cm} every step --- that is what makes the function
\textbf{linear} and the graph \blank{2.0cm}. We call that constant the \textbf{slope}.
\begin{center}
\begin{minipage}[c]{0.46\linewidth}
\[ m \;=\; \frac{\Delta y}{\Delta x} \;=\; \frac{\text{rise}}{\text{run}}
        \;=\; \frac{y_2-y_1}{x_2-x_1} \]
Every linear function can be written
\[ f(x) \;=\; mx + b. \]
\end{minipage}\hfill
\begin{minipage}[c]{0.48\linewidth}\centering
\begin{tikzpicture}[scale=0.44]
  \lgrid{-2}{6}{-6}{6}
  \draw[very thick, forest] (-1,-6)--(5,6);
  \draw[semithick, navy, dashed] (2,0)--(3,0)--(3,2);
  \node[navy, font=\tiny] at (2.5,-0.8){run $1$};
  \node[navy, font=\tiny, right] at (3.1,1){rise $2$};
  \fill[charcoal] (0,-4) circle (3.4pt);
  \fill[charcoal] (2,0) circle (3.4pt);
\end{tikzpicture}\\[1pt]
{\footnotesize $f(x)=2x-4$}
\end{minipage}
\end{center}
\textbf{Slope from any two points.} The line above passes through $(1,-2)$ and $(3,2)$.
\begin{work}
  m &= \frac{2 - (-2)}{3 - 1} \\
    &= \frac{4}{2} \\
    &= 2
\end{work}
\emph{Any} two points give the same answer --- because the rate of change is \blank{2.4cm}.
\end{notesbox}

\boxguard[22]
\begin{notesbox}{2. Where the Graph Meets the Axes}
\begin{center}
\renewcommand{\arraystretch}{1.4}
\begin{tabularx}{\linewidth}{@{}l X X@{}}
\textbf{Feature} & \textbf{How to find it} & \textbf{What it means in a story} \\
\hline
$y$-intercept $(0,b)$ & set $x=\blank{0.7cm}$ & the \blank{2.6cm} value \\
$x$-intercept, or \textbf{zero}, $(a,0)$ & solve $f(x)=\blank{0.7cm}$ & the moment the quantity \blank{2.4cm} \\
\end{tabularx}
\end{center}
\textbf{Worked: $f(x)=2x-4$.} \; Its $y$-intercept is $(0,\blank{0.8cm})$ --- read it straight
off the rule. For the zero, solve:
\begin{work}
  2x - 4 &= 0 \\
      2x &= 4 \\
       x &= 2
\end{work}
so the zero is at $(\blank{0.8cm},\,\blank{0.8cm})$. Check it against the graph above.
\end{notesbox}

\boxguard[26]
\begin{notesbox}{3. Which Way It Moves --- and Where It Sits}
These are \emph{two different questions}. Keep them apart.
\begin{center}
\renewcommand{\arraystretch}{1.35}
\begin{tabularx}{\linewidth}{@{}X|X@{}}
\textbf{Which way it moves} (read left to right) & \textbf{Where it sits} (above or below the $x$-axis) \\
\hline
$m>0$: the line is \blank{2.4cm} & $f(x)>0$ where the graph is \blank{1.8cm} the axis \\
$m<0$: the line is \blank{2.4cm} & $f(x)<0$ where the graph is \blank{1.8cm} the axis \\
$m=0$: the line is \blank{2.4cm} & the sign can only switch at a \blank{1.8cm} \\
\end{tabularx}
\end{center}
\textbf{For $f(x)=2x-4$:} it is \blank{2.2cm}; it is positive when $x \blank{1.2cm}$
and negative when $x \blank{1.2cm}$.

\vspace{3pt}
\begin{tcolorbox}[colback=goldbg, colframe=goldacc, boxrule=0.8pt, arc=1.5mm,
    left=2.5mm, right=2.5mm, top=1.5mm, bottom=1.5mm]
\textbf{\textcolor{goldacc}{Careful: ``positive'' is not ``increasing.''}} \;
Sunday's temperature was $g(h)=6-2h$. At $h=1$ it is $g(1)=\blank{1.0cm}$ degrees --- a
\blank{1.8cm} number --- and the temperature is \blank{1.8cm}. Both at once.
\emph{Positive} says where the graph \textbf{sits}; \emph{increasing} says which way it
\textbf{moves}.
\end{tcolorbox}
\end{notesbox}

\boxguard[20]
\begin{notesbox}{4. Domain and Range --- the Math, and the Story}
\begin{itemize}[leftmargin=1.5em, itemsep=3pt]
  \item A line that is \emph{not} horizontal: domain \blank{2.6cm}, range \blank{2.6cm}.
        The rule accepts every input and produces every output.
  \item A \textbf{horizontal} line $c(x)=b$ is the exception: slope \blank{0.9cm}, neither
        increasing nor decreasing, and range \blank{1.4cm} --- a single value. It has no zero
        at all unless $b=\blank{0.8cm}$, and then \emph{every} input is one.
  \item In a \textbf{story}, the domain shrinks to the inputs that \blank{3.4cm}. Saturday's
        rule works for $h=-100$; Saturday morning does not.
\end{itemize}
\end{notesbox}

\begin{practicebox}
Everything at once, on a new line: $h(x) = -2x + 6$.
\begin{enumerate}[leftmargin=1.7em, itemsep=3pt]
  \item Slope \blank{1.2cm}; \quad $y$-intercept \blank{1.8cm}; \quad increasing or decreasing?
        \blank{2.2cm}
  \item Find the zero.
        \begin{work}
          -2x + 6 &= 0 \\
             -2x  &= -6 \\
               x  &= 3
        \end{work}
        Zero: $(\blank{0.8cm},\,\blank{0.8cm})$.
  \item $h(x)>0$ when $x \blank{1.2cm}$; \quad $h(x)<0$ when $x \blank{1.2cm}$.
  \item Domain \blank{2.2cm}; \quad range \blank{2.2cm}.
\end{enumerate}
\end{practicebox}

\end{document}
